\section{Results: Transformer measurements and falsifiers}
\label{sec:results}

\subsection{RQ1: width and effective data produce a learning response}

\claim{Observation.} Figure~\ref{fig:reference-response}A shows raw outer-seed
task order and 95\% intervals. At $d=16$, the duplicated 32-example controls
produce identical aggregates, after which order increases modestly at 128
examples. Across width, the response is stronger: at $(d,\alpha)=(64,8)$,
semantic order is \ReferenceLargestOrder{} and IID risk is
\ReferenceLargestRisk{}. The resource coordinate, not the nominal control
label, explains the duplicated point.

\claim{Interpretation.} The registered small Transformer exhibits finite-width
learning and benefits from effective data. This validates the task order as a
responsive measurement on the tested grid.

\claim{Not established.} A monotone improvement is not a transition. The grid
has only three widths, one duplicated point, and 24 optimization steps. It does
not identify a thermodynamic scale path or asymptotic threshold.

\claim{Next falsifier.} Use distinct effective sample counts, longer training,
five sizes with the largest held out, and a dense control grid fixed before the
new seeds are observed.

\subsection{Generalization and OOD behavior are separate observables}

At width 64, the low-data generalization gap is \ReferenceLowDataGap{}, whereas
the high-data gap is \ReferenceHighDataGap{}. The OOD-minus-IID differences in
Figure~\ref{fig:reference-response}D are small relative to conservative summed
intervals across this particular synthetic shift.

The result is a data-limited overfitting signature, not an OOD robustness claim.
The OOD generator shares the registered task family, and the five-seed
intervals are too wide to rank small penalties. Reporting train, IID, and OOD
risks separately prevents a reduction in one from being described as a general
improvement in all three.

\manuscriptfigure{figures/figure5_reference_response.pdf}{Public GPU reference
study. Points and lines are condition means with two-sided 95\% Student-$t$
intervals over five independent full-pipeline seeds; translucent points are the
raw seed values. (A) Semantic task order. (B) Normalized IID risk. (C) IID--train
generalization gap. (D) OOD--IID difference with a conservative interval formed
from the two marginal intervals. The nominal control is shown, while
Figure~\ref{fig:reference-verdict} exposes the duplicated effective point.}
{fig:reference-response}

\subsection{RQ2: registered phase diagnostics reject a boundary claim}

\claim{Observation.} The largest-width susceptibility is maximal at the left
endpoint rather than at an interior control. Response peaks do not exhibit a
stable drift sequence. Binder cumulants remain near $2/3$ without a crossing,
and attention entropy remains close to $\log 2$ rather than showing head
specialization (Figure~\ref{fig:reference-mechanisms}B--D).

\claim{Interpretation.} The reference grid supports the outcome
\emph{finite-width learning response; phase boundary not supported}. Its
evidence vector has semantic and outer-seed coordinates, but lacks finite-size
and prospective-prediction support.

\claim{Not established.} Neither criticality, universality, a positional--
semantic transition, nor a Binder crossing is observed. The absence of evidence
on this small grid is not evidence that larger or longer-trained Transformers
cannot transition.

\claim{Next falsifier.} The dense confirmation should test whether an interior
peak emerges after removing effective-control duplication and optimization
censoring. If it does not, the response should remain classified as rounded or
censored.

\subsection{RQ3: module interventions reveal a width-dependent MLP effect}

\claim{Observation.} In the reference grid the largest registered MLP effect is
\ReferenceLargestMLPEffect{}, while the largest attention effect at the
high-data, largest-width point is \ReferenceLargestAttentionEffect{}. MLP
effects grow with width, while attention effects remain smaller on this shallow,
short-trained retrieval task (Figure~\ref{fig:reference-mechanisms}A).

\claim{Interpretation.} The feed-forward path causally contributes to the
registered endpoint under matched ablation. The effect is an $O(1)$ risk
difference with outer-seed uncertainty.

\claim{Not established.} This does not show that MLPs dominate attention in
general, that the modules define separate phases, or that the effect is stable
under depth, long training, or natural language.

\manuscriptfigure{figures/figure6_reference_mechanisms.pdf}{Mechanism and phase
diagnostics from the same reference models. (A) Bounded matched-ablation effects
for attention (dashed) and MLP (solid). (B) Attention entropy with $\log 2$
reference. (C) Outer-seed susceptibility; the largest-width peak is an endpoint.
(D) Binder cumulant with $2/3$ reference. The causal effect in panel A does not
upgrade the negative finite-size evidence in panels B--D.}
{fig:reference-mechanisms}

\manuscriptfigure{figures/figure7_reference_verdict.pdf}{Audit and evidence
verdict. (A) Semantic order and effective training-example count for every
reference condition; the red outline identifies two nominal controls that
collapse to the same 32-example point. (B) Evidence-vector components. The
result is retained as a grade-C finite-width learning response, not converted to
a phase boundary.}{fig:reference-verdict}

\subsection{RQ3--RQ4: frozen confirmations bound stronger conclusions}

The twelve-seed suite tests hypotheses selected from the broad screen. Linear
and GELU blocks reach task orders \ConfirmMLPLinearOrder{} and
\ConfirmMLPGELUOrder{}, compared with \ConfirmMLPGEGLUOrder{} and
\ConfirmMLPSwiGLUOrder{} for matched gated variants. Their bounded MLP effects
are \ConfirmMLPLinearCausalEffect{}, \ConfirmMLPGELUCausalEffect{},
\ConfirmMLPGEGLUCausalEffect{}, and \ConfirmMLPSwiGLUCausalEffect{}. The task
therefore separates these registered architectures, but the causal intervals do
not justify a universal ranking of gating mechanisms.

On frozen learning-rate slices, normalized IID risks are \ConfirmSGDMError{}
(SGD with momentum), \ConfirmAdamWError{} (AdamW), \ConfirmMuonError{} (Muon),
and \ConfirmSOAPError{} (SOAP). These are optimizer-conditioned outcomes inside
the registered grid. They are not a general optimizer benchmark because update
scale, compute, and hyperparameter budgets are not exhaustively matched
\cite{loshchilov2019adamw,tensor_muon_jordan2025,tensor_soap_vyas2024}.

The uncapped data paths produce normalized risks
\ConfirmLinearDataScalingError{}, \ConfirmThreeHalfDataScalingError{}, and
\ConfirmQuadraticDataScalingError{} for $n\propto d$, $d^{3/2}$, and $d^2$ at
the largest common endpoint. Natural-data byte models produce
\ConfirmTinyStoriesBitsPerByte{}, \ConfirmSimpleStoriesBitsPerByte{},
\ConfirmFineWebEduBitsPerByte{}, and \ConfirmDolmaBitsPerByte{}. Differences
across corpora mix modeling difficulty and corpus statistics; they validate a
common bits-per-byte measurement, not a corpus quality ranking.

\manuscriptfigure{figures/figure8_confirmation_summary.pdf}{Unified summary of
the frozen twelve-seed confirmations. Error bars are 95\% Student-$t$ intervals
over fresh full-pipeline seeds. (A) Task order and matched MLP causal effect.
(B) Optimizer-conditioned IID risk on frozen slices. (C) Uncapped data-scaling
paths. (D) Disjoint-range natural-corpus bits per byte. Panel values are separate
estimands and are not pooled into a single score.}{fig:confirmation-summary}

\FloatBarrier
